% Candidate insertion for experimental/grande_finale.tex.
% Task: SP / primitive shift-pair control.
% Status: PROVED SUBCASE / QUOTIENT SIEVE.
%
% This snippet does not prove the full primitive shift-pair bound.  It proves
% a certificate-level quotient-pullback criterion: every shift pair whose two
% locator polynomials have a common nontrivial coefficient scale is exactly a
% quotient pullback, descends with the correct depth, and must be removed from
% the primitive SP ledger.  The remaining obstruction is therefore genuinely
% coefficient-primitive.

\subsection{A quotient sieve for shift pairs}
\label{subsec:sp-next-quotient-sieve}

The exact second-moment ledger leaves a primitive shift-pair term.  Before
trying to bound that term, one needs a checkable quotient sieve: a certificate
should be able to decide whether a putative shift pair is genuinely primitive,
or whether it is a pullback from a smaller quotient scale.  The following
lemmas make that test algebraic.  Throughout this subsection
\[
        D=\alpha H\subseteq \B^\times,\qquad |H|=n,
\]
with \(H\) cyclic and \(\operatorname{char}\B\nmid n\).  For \(c\mid n\),
write \(H_c\le H\) for the subgroup of order \(c\), and
\(\pi_c:D\to D_c\) for \(x\mapsto x^c\).

\begin{definition}[coefficient scale]
\label{def:sp-next-coeff-scale}
Let
\[
        L(X)=X^e+\sum_{j=1}^e \lambda_j X^{e-j}
\]
be a monic polynomial split over \(D\).  Its coefficient scale is
\[
        s(L):=\gcd\bigl(n,\ e,\ \{j:\lambda_j\ne0\}\bigr),
\]
where the gcd of an empty set of indices is interpreted as \(\gcd(n,e)\).
For an ordered pair \((A,B)\) of monic split degree-\(e\) polynomials, put
\[
        s(A,B):=\gcd(s(A),s(B)).
\]
Thus \(s(A,B)>1\) means that both locators have only coefficient gaps
divisible by a common nontrivial divisor of \(n\).
When the same notation is applied on a quotient coset \(D_c\), the ambient
length in the gcd is \(|D_c|=n/c\).
\end{definition}

\begin{lemma}[periodicity, coefficient support, and quotient form]
\label{lem:sp-next-periodicity-equivalence}
Let \(S\subseteq D\), \(|S|=e\), and \(L_S=\ell_S\).  For a divisor
\(c\mid n\), the following are equivalent.
\begin{enumerate}
\item \(S\) is a union of complete \(H_c\)-orbits.
\item \(L_S(X)=L_{S,c}(X^c)\) for a unique monic polynomial
      \(L_{S,c}\) of degree \(e/c\), split over \(D_c=\pi_c(D)\).
\item \(c\mid e\), and in
      \(L_S=X^e+\sum_{j=1}^e\lambda_jX^{e-j}\) one has
      \(\lambda_j=0\) unless \(c\mid j\).
\end{enumerate}
Consequently \(s(L_S)\) is the maximal periodicity scale of \(S\).
\end{lemma}

\begin{proof}
If \(S\) is a union of complete \(H_c\)-orbits, then each orbit over
\(y\in D_c\) has locator \(X^c-y\).  Multiplying over the quotient image
\(\overline S:=\pi_c(S)\) gives
\[
        L_S(X)=\prod_{y\in\overline S}(X^c-y)=L_{S,c}(X^c),
\]
where \(L_{S,c}(Y)=\prod_{y\in\overline S}(Y-y)\).  This proves
\((1)\Rightarrow(2)\), and uniqueness follows from injectivity of
\(L(Y)\mapsto L(X^c)\).

The implication \((2)\Rightarrow(3)\) is immediate by comparing exponents.
Conversely, if the coefficient condition in \((3)\) holds, then
\[
        L_S(X)=X^e+\sum_{c\mid j}\lambda_jX^{e-j}
              =L_{S,c}(X^c)
\]
for a unique monic \(L_{S,c}\in\B[Y]\) of degree \(e/c\).  If
\(y\in D_c\) is a root of \(L_{S,c}\), then every \(x\in D\) with
\(x^c=y\) is a root of \(L_S\); conversely every root of \(L_S\) maps to a
root of \(L_{S,c}\).  Since \(X^n-\alpha^n\) is squarefree, the fibers of
\(\pi_c\) on \(D\) have exactly \(c\) distinct points, so the roots of
\(L_S\) are precisely full \(H_c\)-orbits.  This gives \((3)\Rightarrow(1)\).

The last sentence follows because the set of divisors \(c\mid n\) satisfying
the equivalent conditions is exactly the set of divisors of
\(\gcd(n,e,\{j:\lambda_j\ne0\})\).
\end{proof}

\begin{theorem}[maximal quotient extraction for shift pairs]
\label{thm:sp-next-maximal-extraction}
Let \((A,B)=(\ell_S,\ell_T)\) be a depth-\(w\) shift pair of degree \(e\)
over a multiplicative coset \(D\), meaning \(S,T\subseteq D\) are disjoint
\(e\)-subsets and
\[
        \deg(A-B)\le e-w-1 .
\]
Let \(c=s(A,B)\).  If \(c>1\), then the pair is not primitive: there are
unique monic degree-\(e/c\) locators \(A_c,B_c\), split over \(D_c=\pi_c(D)\),
such that
\[
        A(X)=A_c(X^c),\qquad B(X)=B_c(X^c).
\]
Moreover \((A_c,B_c)\) is a depth
\[
        w_c=\left\lceil {w+1\over c}\right\rceil-1
\]
shift pair on the quotient coset.  If \(c=s(A,B)\) is maximal, then the
descended pair \((A_c,B_c)\) is coefficient-primitive:
\[
        s(A_c,B_c)=1.
\]
Conversely, every coefficient-primitive depth-\(w_c\) quotient pair pulls
back to a depth-\(w\) upstairs pair whose maximal common coefficient scale is
exactly \(c\), provided the two pulled-back root sets are disjoint.
\end{theorem}

\begin{proof}
Since \(c=s(A,B)\), \cref{lem:sp-next-periodicity-equivalence} gives
unique quotient locators \(A_c,B_c\) with \(A=A_c(X^c)\) and \(B=B_c(X^c)\).
The depth condition becomes
\[
        c\deg(A_c-B_c)=\deg(A-B)\le ce_c-w-1,\qquad e_c=e/c.
\]
Equivalently,
\[
        \deg(A_c-B_c)\le e_c-\left\lceil {w+1\over c}\right\rceil,
\]
which is the depth-\(w_c\) condition.

It remains to justify the maximal-scale assertion.  Here coefficient scale for
\((A_c,B_c)\) is computed relative to the quotient length \(n/c\).  Suppose
the quotient pair had common coefficient scale \(d>1\).  By
\cref{lem:sp-next-periodicity-equivalence}, \(A_c=A_{cd}(Y^d)\) and
\(B_c=B_{cd}(Y^d)\) for quotient locators at the further quotient scale.
Then
\[
        A(X)=A_{cd}(X^{cd}),\qquad B(X)=B_{cd}(X^{cd}),
\]
so \(cd\) divides the common coefficient scale of \((A,B)\), contradicting
maximality of \(c\).  Thus \(s(A_c,B_c)=1\).

Conversely, if \(A_c,B_c\) form a depth-\(w_c\) quotient pair, then
\[
        \deg(A_c-B_c)\le e_c-\left\lceil {w+1\over c}\right\rceil .
\]
Multiplying degrees by \(c\) gives
\[
        \deg(A_c(X^c)-B_c(X^c))
        \le ce_c-c\left\lceil {w+1\over c}\right\rceil
        \le ce_c-w-1,
\]
so the pullback is a depth-\(w\) pair upstairs.  Its coefficient scale is at
least \(c\).  If it were larger than \(c\), then after dividing by \(c\) the
extra common divisor would give the quotient pair a nontrivial common
coefficient scale relative to \(n/c\), contrary to the coefficient-primitivity
assumption.
\end{proof}

\begin{corollary}[primitive exclusion by coefficient support]
\label{cor:sp-next-primitive-exclusion}
Any depth-\(w\) shift pair \((A,B)\) with \(s(A,B)>1\) belongs to the
quotient-pullback ledger and must not be counted in the primitive SP cell.
Equivalently, every genuinely primitive shift pair has
\[
        \gcd\bigl(n,\ e,\ \{j:\lambda_j(A)\ne0\},\ \{j:\lambda_j(B)\ne0\}\bigr)=1 .
\]
At the top stratum \(e=w+1\), all binomial prototypes
\[
        A=X^{w+1}-\beta,\qquad B=X^{w+1}-\beta'
\]
with \(A\) and \(B\) split over \(D\) are quotient-borne; they have
coefficient scale \(w+1\) and can occur only when \(w+1\mid n\).
\end{corollary}

\begin{proof}
The first assertion is exactly \cref{thm:sp-next-maximal-extraction}.  The
displayed gcd is the definition of \(s(A,B)\).  For the final assertion,
both binomial locators have only the leading and constant coefficients, so
their common coefficient scale is \(\gcd(n,w+1)\).  If they are split over
\(D=\alpha H\), their root sets are full fibers of \(x\mapsto x^{w+1}\), so
\(w+1\mid n\), and the common coefficient scale is \(w+1\).
\end{proof}

\begin{remark}[remaining primitive obstruction]
\label{rem:sp-next-remaining-obstruction}
\Cref{cor:sp-next-primitive-exclusion} is a genuine narrowing of the SP
input: after this quotient sieve, the primitive ledger contains only pairs
whose two locators have no common coefficient scale.  The theorem does not
bound those coefficient-primitive pairs.  The remaining obstruction is a
family of disjoint split locators \(A,B\), with
\(\deg(A-B)\le e-w-1\), whose coefficient supports have gcd \(1\) after all
common-GCD, tangent, quotient, extension, planted, and bounded-chart cells are
removed.  Proving that such families have only density-model size, or finding
a super-polynomial family of them, is the unresolved SP input.
\end{remark}
